\section{Background}

\subsection{Consumer Device Characteristics}

Consumer devices in our deployment class typically have the following characteristics:
\begin{itemize}
\item \textbf{Resource-Constrained}: Mainstream laptops equipped with 8GB VRAM, 16-32GB RAM, far below data center servers
\item \textbf{Single-Machine Deployment}: No need for distributed coordination, but requires lightweight and efficient systems
\item \textbf{High Stability Requirements}: In home/industrial scenarios, system crashes severely impact user experience
\end{itemize}

\subsection{Multi-Modal Task Characteristics}

We characterize workloads on consumer-grade local AI runtimes using three orthogonal properties: service semantics, device affinity, and latency sensitivity. This view captures local runtime conflicts more directly than a taxonomy based only on model families.

\textbf{Service Semantics}:
\begin{itemize}
\item \textbf{Resident Tasks}: Long-lived services that should remain continuously available and cannot tolerate extended periods without response. Interactive LLM assistants are the canonical example.
\item \textbf{One-shot Tasks}: Bursty requests that are submitted, executed, and completed independently. Image generation is the representative example in this paper.
\end{itemize}

\textbf{Device Affinity}:
\begin{itemize}
\item \textbf{GPU-pinned Tasks}: Tasks whose latency is unacceptable without exclusive GPU execution. In practice, diffusion-based image generation falls into this category on commodity laptops.
\item \textbf{Migratable Tasks}: Tasks that benefit substantially from GPU acceleration but can degrade to CPU execution while remaining usable, depending on quantization and the inference stack. Quantized LLM serving is the representative example.
\end{itemize}

\textbf{Latency Sensitivity}:
\begin{itemize}
\item \textbf{Interactive Tasks}: User-facing requests that require low delay and continuous responsiveness.
\item \textbf{Queueable Tasks}: Tasks for which moderate waiting time is acceptable as long as they eventually complete successfully.
\end{itemize}

Under this abstraction, the current EdgeSched instantiation focuses on a particularly conflict-prone pair in local AI runtimes: LLM serving as a resident, migratable, latency-sensitive service, and image generation as a bursty, GPU-pinned, queueable job. The rest of the paper uses this pair as the concrete instance through which the scheduling policy is presented and evaluated.

\subsection{Design Goals}

Based on the above workload abstraction, EdgeSched's design goals include:

\begin{enumerate}
\item \textbf{Lightweight}: The scheduler itself must remain lightweight (fast startup, small memory), with model memory/VRAM accounted for separately.
\item \textbf{Service Continuity under Resource Degradation}: Resident interactive tasks should preserve availability under GPU contention, even if they must temporarily degrade to CPU execution.
\item \textbf{High Stability}: Crash recovery time should remain on the order of 10 ms, avoiding task state loss and dangling execution states.
\item \textbf{Operational Simplicity}: The system should provide a simple interface and minimal deployment complexity for single-node local AI runtimes.
\end{enumerate}
